\subsection{The Hierarchy Problem and New Dimensions at a Millimeter --- arXiv:hep-ph/9803315}

\textbf{Authors:} Nima Arkani-Hamed, Savas Dimopoulos, Gia Dvali

\paragraph{主な主張}
標準模型を4次元brane上に局在させ重力のみを $n\geq2$ 個の大きなコンパクト余剰次元へ伝播させることで、4次元Planck scaleの大きさを余剰次元の体積効果として説明し、基本的な量子重力スケールを弱スケール近傍まで下げられることを提案した \cite{ArkaniHamed:1998ADD}。

\paragraph{新規性}
弱スケールとPlanck scaleの階層を新たな高エネルギースケールではなく大きな余剰次元の幾何学的体積に置き換える ADD 機構を提示し、sub-millimeter 重力実験や collider で検証可能な phenomenology を初めて具体化した。

\paragraph{位置付け}
large extra dimensions による hierarchy problem 解決の原著論文であり、後続の bulk-neutrino LED 模型や IceCube を用いた KK sterile-neutrino 探索の高次元背景を与える一方、bulk neutrino や $T^2$ complex structure に依存する KK spectrum 自体は扱っていない。
